% Generated by scripts/convert_report.py from source/report.md.


\FloatBarrier

\Needspace{6\baselineskip}

\section{KẾT LUẬN VÀ HƯỚNG PHÁT TRIỂN}


\subsection{Kết quả đạt được và hạn chế}


Đề tài đã xây dựng và kiểm chứng thành công nền tảng MLOps trên AWS, tích hợp các chức năng chính của vòng đời Machine Learning, từ huấn luyện và quản lý mô hình đến triển khai, giám sát và kiểm soát phát hành.


Những kết quả đạt được bao gồm:


\begin{itemize}

\item 

\textbf{Hạ tầng và quản lý tài nguyên:} Sử dụng Terraform để triển khai hạ tầng AWS, kết hợp Amazon EKS và Karpenter để điều phối tài nguyên. Workload huấn luyện được triển khai trên Spot Instances, trong khi dịch vụ suy luận sử dụng On-Demand Instances.


\item 

\textbf{Huấn luyện và quản lý mô hình:} Xây dựng Training Pipeline với LightGBM, tích hợp MLflow để theo dõi thí nghiệm và quản lý phiên bản. Validation Gate hỗ trợ kiểm duyệt chất lượng mô hình trước khi chuyển sang giai đoạn phát hành.


\item 

\textbf{Tự động hóa triển khai:} Tích hợp GitHub Actions, FastAPI, Argo CD và Argo Rollouts để thực hiện CI/CD và Canary Deployment. Hệ thống đã kiểm chứng thành công quy trình tự động toàn trình, bao gồm cả khả năng từ chối bản phát hành không đạt yêu cầu và khôi phục phiên bản Stable.


\item 

\textbf{Giám sát và mở rộng:} Sử dụng Prometheus và Grafana để theo dõi các chỉ số vận hành. HPA cho phép tự động điều chỉnh số lượng replica theo mức sử dụng CPU.


\end{itemize}

Bên cạnh những kết quả trên, hệ thống vẫn còn một số hạn chế. Kiến trúc hiện tại chưa đáp ứng đầy đủ yêu cầu \textbf{High Availability} và \textbf{Disaster Recovery}. Mặc dù đã có cơ chế backup/restore, khả năng khôi phục khi xảy ra sự cố lớn vẫn cần được đánh giá thêm.


Đối với quá trình vận hành, cơ chế Canary và rollback mới được kiểm chứng trong các kịch bản xác định, chưa bao phủ mọi tình huống lỗi. HPA đã thể hiện khả năng tự động mở rộng nhưng chưa có thử nghiệm đối chứng để đánh giá hiệu quả cải thiện thông lượng. Ngoài ra, dữ liệu giám sát chưa được lưu trữ dài hạn và hệ thống chưa hỗ trợ phát hiện Data Drift để tự động kích hoạt tái huấn luyện.


\subsection{Hướng phát triển}


Từ các kết quả và hạn chế trên, đề tài có thể tiếp tục được phát triển theo những hướng chính:


\begin{itemize}

\item 

\textbf{Nâng cao tính sẵn sàng và khả năng phục hồi:} Hoàn thiện kiến trúc High Availability, triển khai các thành phần quan trọng trên nhiều Availability Zones và xây dựng chiến lược Disaster Recovery với các chỉ tiêu RPO, RTO cụ thể.


\item 

\textbf{Cải thiện giám sát và độ tin cậy khi phát hành:} Bổ sung lưu trữ metrics dài hạn, cảnh báo tự động và kiểm thử thêm các tình huống lỗi Canary. Đồng thời, thực hiện thử nghiệm tải có đối chứng để đánh giá hiệu quả mở rộng của HPA và Karpenter.


\item 

\textbf{Hoàn thiện vòng đời MLOps tự động:} Triển khai Data Drift Detection để phát hiện sự thay đổi của dữ liệu, từ đó kích hoạt tái huấn luyện và đánh giá mô hình. Các phiên bản mới tiếp tục được kiểm soát thông qua Validation Gate và Canary Deployment trước khi phát hành.


\item 

\textbf{Tăng cường bảo mật:} Hoàn thiện quản lý secrets, phân quyền truy cập và kiểm tra an toàn cho container images cùng các quy trình CI/CD.


\end{itemize}

Nhìn chung, đề tài đã hoàn thành mục tiêu xây dựng và kiểm chứng một nền tảng MLOps trên AWS trong phạm vi thực nghiệm, tạo cơ sở để tiếp tục phát triển thành hệ thống có độ tin cậy và khả năng tự động hóa cao hơn.

